\songtitle{La sombra y el futuro}

\tag*{2022}\tag{Scholer}\tag{clan-scholer}
\begin{notes}
Lyrics by Claude and Carlos Thompson, based on characters, %
from the character dossier document for the planed series \emph{Scholer} %
written and conceptualized by Carlos Thompson.
\end{notes}
\begin{style}
latin romantic ballad, 1980s soft rock, orchestral pop, electric piano arpeggios, nylon guitar strums, sustained lead guitar, saxophone counterlines, string ensemble swells, analog synth pads, tom fills, gang vocals, stacked harmonies, plate reverb, tape saturation, gated snare, wide stereo mix, bittersweet nostalgia, anthemic chorus, 72 BPM, half-time groove
\end{style}
\begin{lyrics}
\songpart{Verse 1}
Crecen los cuatro bajo un mismo techo,\\
con la ciudad entera mirando su futuro,\\
hijos del apellido que en Bogotá se asoma,\\
la nueva generación que crece sin apuro.

\songpart{Verse 2}
Agustín construye planos en la ingeniería,\\
María Luisa regresa, con Derecho en la mira,\\
Andrea Sofía se gradúa en unos meses,\\
y Nicolás, el menor, apenas empieza su vida.

\songpart{Chorus}
Se abre el futuro, ancho y sin condena,\\
universidades, planes, un mundo que empieza,\\
pero detrás camina una sombra ajena,\\
la historia de la casa que nadie les enseña.

\songpart{Verse 3}
Andrea Sofía tiene una amiga entrañable,\\
Sandra, su mejor amiga, un lazo inseparable,\\
no saben, las dos, que su cercanía repite\\
un vínculo mucho más viejo, secreto y notable.

\songpart{Bridge}
Llevan un apellido que pesa más de lo que creen,\\
una casa de silencios que ellos no escogieron,\\
crecieron entre lujos y entre pactos que no ven,\\
y un día, tal vez, entenderán lo que vivieron.

\songpart{Final Chorus}
Se abre el futuro, ancho y sin condena,\\
universidades, planes, un mundo que empieza,\\
pero detrás camina una sombra ajena,\\
la historia de la casa que nadie les enseña.\\
Algún día el pasado los alcanza y despierta,\\
y entonces sabrán que su casa tenía otra puerta.
\end{lyrics}
